% !TEX root = ../main.tex
%% Corrigé de l'examen de fin d'année — vérifié par verifications/eval_ex.py.

\corrigeSujetEval{EX2}{Examen de fin d'année}

\begin{corrigeEval}{4 points}{Points fixes d'une permutation aléatoire}
\begin{questions}
  \item Les permutations fixant $i$ correspondent aux permutations des
    $n-1$ autres éléments : $P(X_i=1)=\frac{(n-1)!}{n!}=\frac1n$. Comme
    $X=\sum_iX_i$, par linéarité $\res{E(X)=n\cdot\frac1n=1}$. \pts{1 pt}
  \item De même $P(X_i=1,X_j=1)=\frac{(n-2)!}{n!}=\frac{1}{n(n-1)}$. Avec
    $X_i^{2}=X_i$ :
    \[E(X^{2})=\sum_{i}E(X_i)+\sum_{i\neq j}E(X_iX_j)=1+n(n-1)\cdot\frac{1}{n(n-1)}=2,\]
    d'où $\res{V(X)=2-1^{2}=1}$ : moyenne et variance ne dépendent pas de $n$.
    \pts{1,5 pt}
  \item Seule l'identité a $n$ points fixes : $\res{P(X=n)=\frac{1}{n!}}$. Une
    permutation qui fixe $n-1$ points fixe aussi le dernier :
    $\res{P(X=n-1)=0}$. Enfin
    $P(X\geqslant3)\leqslant P\bigl(|X-1|\geqslant2\bigr)\leqslant\frac{V(X)}{2^{2}}=\res{\frac14}$.
    \pts{1 pt}
  \item On compte les permutations de $S_4$ selon leurs points fixes (les
    dérangements de $k$ éléments sont $1,0,1,2,9$ pour $k=0,\dots,4$) :
    $9$ sans point fixe, $4\times2=8$ avec un seul, $\binom42\times1=6$ avec
    deux, $0$ avec trois, $1$ avec quatre. D'où la loi (figure) et
    $E(X)=\frac{8+12+4}{24}=1$, $E(X^{2})=\frac{8+24+16}{24}=2$.
    \pts{0,5 pt}
\end{questions}
\begin{center}\begin{minipage}{\linewidth}\centering
\begin{tikzpicture}[x=1.15cm,y=11cm,line cap=round]
  \draw[otMuted,line width=0.5pt,-{Stealth[length=1.5mm]}] (-0.6,0) -- (4.7,0);
  \draw[otMuted,line width=0.5pt,-{Stealth[length=1.5mm]}] (-0.6,0) -- (-0.6,0.44);
  \foreach \k/\p/\lab in {0/0.375/{9/24},1/0.3333/{8/24},2/0.25/{6/24},3/0/{0},4/0.04167/{1/24}}{
    \fill[otNavy,rounded corners=0.6pt] ({\k-0.28},0) rectangle ({\k+0.28},\p);
    \node[font=\scriptsize,anchor=south] at (\k,\p) {$\lab$};
    \node[font=\scriptsize,anchor=north] at (\k,-0.005) {$\k$};}
  \node[font=\scriptsize,anchor=west] at (4.7,0) {$k$};
  \node[font=\scriptsize,anchor=south] at (-0.6,0.44) {$P(X=k)$};
\end{tikzpicture}\\*[0.5mm]
\legende{La loi du nombre de points fixes pour $n=4$ : on vérifie
$E(X)=V(X)=1$.}
\end{minipage}\end{center}
\end{corrigeEval}

\begin{corrigeEval}{5 points}{Matrices de rang 1}
\begin{questions}
  \item La $j$-ième colonne de $M$ est $v_jU$ : toutes les colonnes sont
    colinéaires à $U$, et l'une au moins est non nulle ($V\neq0$), donc
    $\res{\operatorname{rg}M=1}$. Puis
    $\operatorname{tr}M=\sum_iu_iv_i=V^{\mathsf{T}}U$, et par associativité
    \[M^{2}=U\,(V^{\mathsf{T}}U)\,V^{\mathsf{T}}=(V^{\mathsf{T}}U)\,UV^{\mathsf{T}}=\res{(\operatorname{tr}M)\,M},\]
    car $V^{\mathsf{T}}U$ est un scalaire. \pts{1,5 pt}
  \item Si $t=\operatorname{tr}M\neq0$, $P=\frac1tM$ vérifie
    $P^{2}=\frac{1}{t^{2}}\,tM=P$ : \res{\text{projecteur}} (de rang $1$). Si
    $t=0$, $M^{2}=0$ : $M$ est \res{\text{nilpotente}}. \pts{1 pt}
  \item Si $t\neq-1$ :
    $(I_n+M)\bigl(I_n-\frac{1}{1+t}M\bigr)=I_n+M-\frac{1}{1+t}(M+tM)=I_n$,
    d'où l'inversibilité et
    $\res{(I_n+M)^{-1}=I_n-\dfrac{1}{1+\operatorname{tr}M}\,M}$. Si $t=-1$ :
    $(I_n+M)U=U+U\,(V^{\mathsf{T}}U)=U-U=0$ avec $U\neq0$ : $I_n+M$ n'est
    \res{\text{pas inversible}}. \pts{1 pt}
  \item Si $t\neq0$, $\frac1tM$ est un projecteur de rang $1$ : dans une base
    adaptée, $M$ est semblable à $\operatorname{diag}(t,0,\dots,0)$, donc
    $I_n+M$ à $\operatorname{diag}(1+t,1,\dots,1)$. Si $t=0$, $M^{2}=0$ donne
    $U\in\operatorname{Im}M\subseteq\operatorname{Ker}M$, de dimension $n-1$ :
    on complète $e_1=U$ en une base $(e_1,\dots,e_{n-1})$ de
    $\operatorname{Ker}M$, puis en une base $(e_1,\dots,e_n)$ de $\R^{n}$.
    Comme $Me_n\in\operatorname{Vect}(e_1)$, $M$ est semblable à une matrice
    dont seul le coefficient $(1,n)$ peut être non nul, et $I_n+M$ à une
    matrice triangulaire de diagonale $1$. Dans les deux cas
    $\res{\det(I_n+M)=1+\operatorname{tr}M}$. \pts{1 pt}
  \item $M=\begin{pmatrix}1&1&1\\2&2&2\\3&3&3\end{pmatrix}$, $\operatorname{tr}M=6$ :
    \[\res{(I_3+M)^{-1}=I_3-\frac17M=\frac17\begin{pmatrix}6&-1&-1\\-2&5&-2\\-3&-3&4\end{pmatrix}},
      \qquad\res{\det(I_3+M)=7}.\]
    \pts{0,5 pt}
\end{questions}
\end{corrigeEval}

\begin{corrigeEval}{4 points}{Le folium de Descartes}
\begin{questions}
  \item $\partial_1f=3x^{2}-3y$ et $\partial_2f=3y^{2}-3x$ : $y=x^{2}$ et
    $x=y^{2}$, d'où $x=x^{4}$, soit $x\in\ens{0,1}$ (racines réelles) :
    les points critiques sont $\res{(0,0)\text{ et }(1,1)}$. \pts{1 pt}
  \item $f(0,0)=0$ et $f(x,0)=x^{3}$ prend des valeurs strictement positives
    et strictement négatives au voisinage de $0$ : \res{\text{pas d'extremum
    local}} en $(0,0)$ (point col). \pts{0,5 pt}
  \item Le développement est un calcul direct. Puis
    $h^{2}-hk+k^{2}-\frac12(h^{2}+k^{2})=\frac12(h-k)^{2}\geqslant0$. Avec
    $r=\sqrt{h^{2}+k^{2}}$, on a $|h|,|k|\leqslant r$, donc
    $|h^{3}+k^{3}|\leqslant|h|^{3}+|k|^{3}\leqslant r\,(h^{2}+k^{2})=r^{3}$, d'où
    \[f(1+h,1+k)+1\geqslant\frac32r^{2}-r^{3}=r^{2}\Bigl(\frac32-r\Bigr)\geqslant0\quad\text{pour }r\leqslant\frac32.\]
    $f$ admet donc en $(1,1)$ un \res{\text{minimum local}, f(1,1)=-1}, qui
    n'est \res{\text{pas global}} : $f(x,0)=x^{3}\to-\infty$ en $-\infty$.
    \pts{2 pts}
  \item Avec $x=\frac{3t}{1+t^{3}}$, $y=tx$ :
    $x^{3}+y^{3}=x^{3}(1+t^{3})=\frac{27t^{3}}{(1+t^{3})^{2}}$ et
    $3xy=3tx^{2}=\frac{27t^{3}}{(1+t^{3})^{2}}$ : \res{f(x,y)=0}. \pts{0,5 pt}
\end{questions}
\begin{center}\begin{minipage}{\linewidth}\centering
\begin{tikzpicture}[scale=1.05,line cap=round,line join=round]
  \begin{scope}
    \clip (-2,-2) rectangle (2.5,2.5);
    % Fichier généré par verifications/fig_folium.py — ne pas éditer
\draw[otGreenDark!80,line width=0.6pt] (-2.000,1.022) -- (-1.984,1.006) -- (-1.968,0.991) -- (-1.951,0.974) -- (-1.936,0.959) -- (-1.918,0.942) -- (-1.903,0.928) -- (-1.885,0.910) -- (-1.871,0.896) -- (-1.852,0.878) -- (-1.839,0.865) -- (-1.820,0.846) -- (-1.806,0.833) -- (-1.788,0.815) -- (-1.773,0.800) -- (-1.755,0.784) -- (-1.739,0.768) -- (-1.723,0.753) -- (-1.706,0.736) -- (-1.691,0.721) -- (-1.673,0.704) -- (-1.659,0.690) -- (-1.639,0.672) -- (-1.626,0.659) -- (-1.607,0.641) -- (-1.593,0.627) -- (-1.575,0.610) -- (-1.559,0.594) -- (-1.543,0.579) -- (-1.526,0.562) -- (-1.511,0.548) -- (-1.492,0.530) -- (-1.479,0.517) -- (-1.459,0.499) -- (-1.445,0.485) -- (-1.427,0.468) -- (-1.411,0.453) -- (-1.395,0.437) -- (-1.377,0.421) -- (-1.363,0.407) -- (-1.343,0.388) -- (-1.330,0.376) -- (-1.311,0.358) -- (-1.296,0.343) -- (-1.279,0.328) -- (-1.262,0.311) -- (-1.247,0.297) -- (-1.227,0.279) -- (-1.214,0.266) -- (-1.195,0.249) -- (-1.179,0.234) -- (-1.163,0.219) -- (-1.145,0.202) -- (-1.131,0.189) -- (-1.112,0.170) -- (-1.097,0.157) -- (-1.079,0.140) -- (-1.062,0.124) -- (-1.047,0.110) -- (-1.028,0.093) -- (-1.014,0.079) -- (-0.996,0.063) -- (-0.979,0.047) -- (-0.964,0.033) -- (-0.944,0.015) -- (-0.931,0.002) -- (-0.912,-0.015) -- (-0.896,-0.030) -- (-0.880,-0.045) -- (-0.861,-0.062) -- (-0.847,-0.075) -- (-0.828,-0.093) -- (-0.813,-0.107) -- (-0.796,-0.123) -- (-0.778,-0.139) -- (-0.764,-0.153) -- (-0.745,-0.171) -- (-0.730,-0.185) -- (-0.712,-0.201) -- (-0.696,-0.217) -- (-0.680,-0.232) -- (-0.662,-0.249) -- (-0.648,-0.262) -- (-0.629,-0.281) -- (-0.615,-0.294) -- (-0.597,-0.311) -- (-0.581,-0.326) -- (-0.564,-0.343) -- (-0.548,-0.358) -- (-0.532,-0.374) -- (-0.515,-0.391) -- (-0.500,-0.406) -- (-0.483,-0.423) -- (-0.468,-0.437) -- (-0.450,-0.455) -- (-0.436,-0.470) -- (-0.418,-0.487) -- (-0.403,-0.502) -- (-0.387,-0.519) -- (-0.371,-0.535) -- (-0.355,-0.552) -- (-0.339,-0.568) -- (-0.324,-0.584) -- (-0.307,-0.601) -- (-0.293,-0.616) -- (-0.275,-0.635) -- (-0.262,-0.648) -- (-0.244,-0.667) -- (-0.230,-0.682) -- (-0.213,-0.700) -- (-0.197,-0.717) -- (-0.183,-0.732) -- (-0.165,-0.751) -- (-0.152,-0.765) -- (-0.135,-0.783) -- (-0.120,-0.799) -- (-0.105,-0.815) -- (-0.088,-0.834) -- (-0.075,-0.848) -- (-0.057,-0.867) -- (-0.043,-0.882) -- (-0.027,-0.899) -- (-0.011,-0.917) -- (0.003,-0.931) -- (0.021,-0.951) -- (0.034,-0.965) -- (0.051,-0.983) -- (0.067,-1.000) -- (0.081,-1.015) -- (0.098,-1.034) -- (0.112,-1.048) -- (0.128,-1.067) -- (0.144,-1.083) -- (0.158,-1.099) -- (0.176,-1.117) -- (0.189,-1.131) -- (0.207,-1.150) -- (0.221,-1.166) -- (0.237,-1.182) -- (0.253,-1.200) -- (0.267,-1.215) -- (0.285,-1.234) -- (0.298,-1.248) -- (0.315,-1.266) -- (0.330,-1.282) -- (0.346,-1.298) -- (0.363,-1.316) -- (0.376,-1.330) -- (0.395,-1.350) -- (0.408,-1.364) -- (0.425,-1.382) -- (0.440,-1.398) -- (0.456,-1.414) -- (0.472,-1.431) -- (0.486,-1.446) -- (0.504,-1.465) -- (0.517,-1.479) -- (0.536,-1.498) -- (0.549,-1.512) -- (0.566,-1.530) -- (0.582,-1.546) -- (0.597,-1.562) -- (0.614,-1.579) -- (0.628,-1.594) -- (0.646,-1.613) -- (0.659,-1.627) -- (0.678,-1.646) -- (0.691,-1.659) -- (0.709,-1.678) -- (0.723,-1.693) -- (0.740,-1.710) -- (0.755,-1.726) -- (0.771,-1.742) -- (0.788,-1.759) -- (0.802,-1.775) -- (0.820,-1.792) -- (0.834,-1.807) -- (0.852,-1.826) -- (0.865,-1.839) -- (0.884,-1.858) -- (0.897,-1.872) -- (0.915,-1.891) -- (0.929,-1.905) -- (0.947,-1.923) -- (0.961,-1.938) -- (0.978,-1.955) -- (0.994,-1.971) -- (1.009,-1.987);
\draw[otGreenDark!80,line width=0.6pt] (0.781,0.677) -- (0.757,0.685) -- (0.736,0.695) -- (0.717,0.709) -- (0.704,0.723) -- (0.691,0.743) -- (0.681,0.768) -- (0.675,0.794) -- (0.674,0.826) -- (0.677,0.858) -- (0.682,0.884) -- (0.689,0.910) -- (0.697,0.934) -- (0.706,0.955) -- (0.717,0.977) -- (0.729,1.000) -- (0.741,1.019) -- (0.754,1.039) -- (0.768,1.058) -- (0.781,1.075) -- (0.794,1.091) -- (0.811,1.109) -- (0.826,1.125) -- (0.843,1.142) -- (0.858,1.156) -- (0.878,1.172) -- (0.897,1.187) -- (0.914,1.200) -- (0.934,1.212) -- (0.955,1.225) -- (0.974,1.236) -- (0.994,1.245) -- (1.019,1.256) -- (1.041,1.264) -- (1.064,1.271) -- (1.096,1.277) -- (1.122,1.279) -- (1.155,1.278) -- (1.180,1.273) -- (1.205,1.264) -- (1.225,1.252) -- (1.241,1.238) -- (1.256,1.219) -- (1.267,1.200) -- (1.275,1.174) -- (1.279,1.148) -- (1.279,1.116) -- (1.276,1.090) -- (1.270,1.063) -- (1.263,1.039) -- (1.254,1.013) -- (1.245,0.992) -- (1.232,0.968) -- (1.221,0.948) -- (1.209,0.929) -- (1.196,0.910) -- (1.182,0.891) -- (1.167,0.872) -- (1.155,0.857) -- (1.138,0.839) -- (1.122,0.823) -- (1.105,0.807) -- (1.090,0.794) -- (1.071,0.778) -- (1.052,0.763) -- (1.032,0.749) -- (1.013,0.737) -- (0.994,0.726) -- (0.974,0.715) -- (0.950,0.704) -- (0.929,0.695) -- (0.903,0.687) -- (0.878,0.680) -- (0.852,0.676) -- (0.820,0.674) -- (0.788,0.676) -- cycle;
\draw[otGreenDark!55,line width=0.6pt] (-2.000,1.060) -- (-1.984,1.045) -- (-1.968,1.030) -- (-1.950,1.013) -- (-1.936,0.999) -- (-1.916,0.981) -- (-1.902,0.968) -- (-1.884,0.951) -- (-1.867,0.936) -- (-1.852,0.921) -- (-1.833,0.903) -- (-1.819,0.891) -- (-1.800,0.874) -- (-1.784,0.858) -- (-1.768,0.844) -- (-1.749,0.826) -- (-1.735,0.813) -- (-1.717,0.797) -- (-1.700,0.781) -- (-1.685,0.767) -- (-1.665,0.750) -- (-1.650,0.736) -- (-1.633,0.721) -- (-1.614,0.704) -- (-1.600,0.691) -- (-1.582,0.674) -- (-1.564,0.659) -- (-1.549,0.645) -- (-1.530,0.628) -- (-1.514,0.614) -- (-1.498,0.600) -- (-1.479,0.583) -- (-1.463,0.569) -- (-1.446,0.554) -- (-1.427,0.537) -- (-1.411,0.524) -- (-1.395,0.509) -- (-1.376,0.493) -- (-1.359,0.479) -- (-1.343,0.465) -- (-1.324,0.448) -- (-1.307,0.433) -- (-1.291,0.421) -- (-1.273,0.405) -- (-1.253,0.388) -- (-1.238,0.376) -- (-1.221,0.361) -- (-1.202,0.345) -- (-1.184,0.330) -- (-1.168,0.318) -- (-1.150,0.303) -- (-1.131,0.287) -- (-1.113,0.273) -- (-1.098,0.260) -- (-1.079,0.245) -- (-1.060,0.229) -- (-1.042,0.215) -- (-1.026,0.202) -- (-1.009,0.188) -- (-0.989,0.173) -- (-0.970,0.158) -- (-0.952,0.144) -- (-0.936,0.131) -- (-0.918,0.117) -- (-0.899,0.102) -- (-0.880,0.087) -- (-0.861,0.073) -- (-0.845,0.060) -- (-0.828,0.047) -- (-0.809,0.033) -- (-0.790,0.018) -- (-0.770,0.003) -- (-0.752,-0.011) -- (-0.736,-0.024) -- (-0.719,-0.036) -- (-0.700,-0.051) -- (-0.680,-0.066) -- (-0.661,-0.081) -- (-0.644,-0.094) -- (-0.627,-0.107) -- (-0.609,-0.121) -- (-0.590,-0.137) -- (-0.571,-0.152) -- (-0.555,-0.165) -- (-0.539,-0.178) -- (-0.519,-0.194) -- (-0.500,-0.210) -- (-0.485,-0.223) -- (-0.468,-0.238) -- (-0.448,-0.255) -- (-0.434,-0.268) -- (-0.416,-0.284) -- (-0.399,-0.300) -- (-0.384,-0.314) -- (-0.365,-0.333) -- (-0.352,-0.346) -- (-0.333,-0.365) -- (-0.320,-0.379) -- (-0.302,-0.397) -- (-0.288,-0.413) -- (-0.272,-0.429) -- (-0.255,-0.448) -- (-0.242,-0.463) -- (-0.227,-0.481) -- (-0.210,-0.500) -- (-0.197,-0.515) -- (-0.184,-0.532) -- (-0.168,-0.552) -- (-0.152,-0.571) -- (-0.139,-0.587) -- (-0.126,-0.603) -- (-0.111,-0.622) -- (-0.096,-0.642) -- (-0.082,-0.660) -- (-0.069,-0.677) -- (-0.056,-0.694) -- (-0.041,-0.712) -- (-0.026,-0.732) -- (-0.012,-0.751) -- (0.002,-0.769) -- (0.015,-0.786) -- (0.028,-0.803) -- (0.043,-0.822) -- (0.057,-0.841) -- (0.072,-0.861) -- (0.086,-0.878) -- (0.099,-0.895) -- (0.112,-0.912) -- (0.127,-0.931) -- (0.142,-0.951) -- (0.157,-0.969) -- (0.170,-0.985) -- (0.183,-1.002) -- (0.198,-1.021) -- (0.214,-1.041) -- (0.227,-1.058) -- (0.240,-1.074) -- (0.255,-1.092) -- (0.271,-1.112) -- (0.285,-1.129) -- (0.298,-1.145) -- (0.313,-1.163) -- (0.329,-1.182) -- (0.343,-1.200) -- (0.356,-1.215) -- (0.372,-1.234) -- (0.388,-1.253) -- (0.401,-1.269) -- (0.415,-1.285) -- (0.432,-1.305) -- (0.446,-1.322) -- (0.459,-1.337) -- (0.476,-1.356) -- (0.491,-1.374) -- (0.504,-1.389) -- (0.521,-1.408) -- (0.536,-1.426) -- (0.549,-1.441) -- (0.566,-1.459) -- (0.582,-1.477) -- (0.594,-1.492) -- (0.611,-1.511) -- (0.627,-1.528) -- (0.640,-1.543) -- (0.657,-1.562) -- (0.672,-1.579) -- (0.686,-1.594) -- (0.703,-1.614) -- (0.717,-1.629) -- (0.732,-1.646) -- (0.749,-1.664) -- (0.762,-1.678) -- (0.779,-1.697) -- (0.794,-1.714) -- (0.809,-1.730) -- (0.826,-1.749) -- (0.839,-1.763) -- (0.856,-1.781) -- (0.871,-1.798) -- (0.886,-1.813) -- (0.903,-1.833) -- (0.916,-1.846) -- (0.933,-1.865) -- (0.948,-1.881) -- (0.963,-1.897) -- (0.981,-1.916) -- (0.994,-1.929) -- (1.012,-1.948) -- (1.026,-1.964) -- (1.042,-1.981) -- (1.058,-1.998);
\draw[otGreenDark!55,line width=0.6pt] (0.601,0.408) -- (0.575,0.411) -- (0.549,0.418) -- (0.525,0.427) -- (0.504,0.438) -- (0.485,0.451) -- (0.468,0.466) -- (0.453,0.482) -- (0.440,0.500) -- (0.428,0.524) -- (0.420,0.543) -- (0.413,0.569) -- (0.408,0.599) -- (0.406,0.627) -- (0.406,0.659) -- (0.409,0.685) -- (0.413,0.717) -- (0.419,0.742) -- (0.425,0.768) -- (0.433,0.794) -- (0.440,0.814) -- (0.450,0.839) -- (0.459,0.862) -- (0.469,0.884) -- (0.479,0.903) -- (0.491,0.928) -- (0.503,0.948) -- (0.514,0.968) -- (0.526,0.987) -- (0.538,1.006) -- (0.551,1.026) -- (0.565,1.045) -- (0.580,1.064) -- (0.594,1.084) -- (0.607,1.099) -- (0.621,1.116) -- (0.638,1.135) -- (0.652,1.151) -- (0.668,1.167) -- (0.685,1.184) -- (0.700,1.200) -- (0.717,1.215) -- (0.735,1.232) -- (0.750,1.245) -- (0.768,1.260) -- (0.788,1.275) -- (0.806,1.290) -- (0.824,1.303) -- (0.842,1.315) -- (0.862,1.328) -- (0.882,1.341) -- (0.903,1.354) -- (0.923,1.365) -- (0.942,1.375) -- (0.965,1.386) -- (0.987,1.397) -- (1.009,1.406) -- (1.032,1.415) -- (1.058,1.423) -- (1.084,1.431) -- (1.109,1.437) -- (1.135,1.442) -- (1.161,1.445) -- (1.193,1.447) -- (1.225,1.447) -- (1.251,1.444) -- (1.283,1.438) -- (1.304,1.431) -- (1.328,1.421) -- (1.348,1.410) -- (1.367,1.396) -- (1.385,1.380) -- (1.399,1.363) -- (1.412,1.344) -- (1.424,1.322) -- (1.432,1.303) -- (1.439,1.277) -- (1.444,1.251) -- (1.447,1.219) -- (1.447,1.187) -- (1.445,1.155) -- (1.441,1.129) -- (1.436,1.103) -- (1.429,1.077) -- (1.421,1.052) -- (1.412,1.026) -- (1.405,1.006) -- (1.394,0.981) -- (1.385,0.961) -- (1.373,0.939) -- (1.361,0.916) -- (1.350,0.897) -- (1.339,0.878) -- (1.326,0.858) -- (1.313,0.839) -- (1.300,0.820) -- (1.285,0.800) -- (1.270,0.781) -- (1.258,0.765) -- (1.244,0.749) -- (1.227,0.730) -- (1.212,0.714) -- (1.197,0.697) -- (1.180,0.680) -- (1.165,0.665) -- (1.148,0.650) -- (1.130,0.633) -- (1.115,0.620) -- (1.097,0.605) -- (1.077,0.589) -- (1.059,0.575) -- (1.041,0.562) -- (1.023,0.549) -- (1.004,0.536) -- (0.984,0.524) -- (0.962,0.511) -- (0.942,0.499) -- (0.923,0.488) -- (0.903,0.479) -- (0.878,0.466) -- (0.858,0.458) -- (0.833,0.447) -- (0.813,0.440) -- (0.788,0.431) -- (0.762,0.424) -- (0.736,0.417) -- (0.710,0.412) -- (0.678,0.408) -- (0.652,0.406) -- (0.620,0.406) -- cycle;
\draw[otNavy!45,line width=0.6pt] (-2.000,1.178) -- (-1.981,1.162) -- (-1.964,1.148) -- (-1.948,1.135) -- (-1.929,1.119) -- (-1.910,1.104) -- (-1.893,1.090) -- (-1.877,1.077) -- (-1.858,1.062) -- (-1.839,1.047) -- (-1.821,1.032) -- (-1.805,1.019) -- (-1.788,1.006) -- (-1.768,0.991) -- (-1.749,0.976) -- (-1.730,0.961) -- (-1.713,0.948) -- (-1.696,0.936) -- (-1.678,0.922) -- (-1.659,0.908) -- (-1.639,0.893) -- (-1.620,0.879) -- (-1.601,0.865) -- (-1.582,0.852) -- (-1.564,0.839) -- (-1.546,0.826) -- (-1.528,0.813) -- (-1.509,0.800) -- (-1.490,0.788) -- (-1.471,0.775) -- (-1.452,0.762) -- (-1.432,0.749) -- (-1.412,0.736) -- (-1.392,0.723) -- (-1.372,0.710) -- (-1.351,0.697) -- (-1.330,0.685) -- (-1.311,0.673) -- (-1.292,0.662) -- (-1.273,0.651) -- (-1.252,0.639) -- (-1.229,0.627) -- (-1.208,0.615) -- (-1.189,0.605) -- (-1.168,0.594) -- (-1.144,0.582) -- (-1.124,0.573) -- (-1.101,0.562) -- (-1.079,0.552) -- (-1.057,0.543) -- (-1.034,0.533) -- (-1.010,0.524) -- (-0.989,0.516) -- (-0.964,0.507) -- (-0.938,0.498) -- (-0.916,0.491) -- (-0.893,0.485) -- (-0.867,0.478) -- (-0.841,0.472) -- (-0.812,0.466) -- (-0.783,0.461) -- (-0.758,0.457) -- (-0.725,0.454) -- (-0.700,0.453) -- (-0.667,0.452) -- (-0.640,0.453) -- (-0.609,0.455) -- (-0.578,0.459) -- (-0.552,0.464) -- (-0.526,0.470) -- (-0.500,0.477) -- (-0.478,0.485) -- (-0.455,0.494) -- (-0.431,0.504) -- (-0.410,0.515) -- (-0.391,0.525) -- (-0.371,0.536) -- (-0.350,0.549) -- (-0.331,0.562) -- (-0.313,0.575) -- (-0.294,0.590) -- (-0.275,0.605) -- (-0.257,0.620) -- (-0.242,0.633) -- (-0.223,0.651) -- (-0.208,0.665) -- (-0.191,0.681) -- (-0.175,0.697) -- (-0.159,0.714) -- (-0.144,0.730) -- (-0.127,0.748) -- (-0.113,0.762) -- (-0.096,0.781) -- (-0.082,0.797) -- (-0.067,0.813) -- (-0.049,0.832) -- (-0.036,0.847) -- (-0.021,0.865) -- (-0.004,0.883) -- (0.009,0.898) -- (0.025,0.916) -- (0.041,0.934) -- (0.054,0.948) -- (0.071,0.968) -- (0.086,0.984) -- (0.100,1.000) -- (0.118,1.019) -- (0.131,1.033) -- (0.148,1.052) -- (0.163,1.068) -- (0.178,1.084) -- (0.195,1.102) -- (0.209,1.116) -- (0.227,1.135) -- (0.240,1.148) -- (0.259,1.167) -- (0.273,1.180) -- (0.292,1.199) -- (0.305,1.212) -- (0.324,1.230) -- (0.339,1.245) -- (0.356,1.260) -- (0.374,1.277) -- (0.389,1.290) -- (0.408,1.307) -- (0.425,1.322) -- (0.440,1.335) -- (0.459,1.351) -- (0.479,1.367) -- (0.495,1.380) -- (0.511,1.393) -- (0.530,1.408) -- (0.549,1.423) -- (0.569,1.437) -- (0.587,1.451) -- (0.605,1.464) -- (0.623,1.476) -- (0.642,1.489) -- (0.661,1.502) -- (0.680,1.515) -- (0.700,1.528) -- (0.721,1.541) -- (0.742,1.554) -- (0.762,1.565) -- (0.781,1.576) -- (0.800,1.587) -- (0.822,1.599) -- (0.845,1.611) -- (0.865,1.621) -- (0.885,1.631) -- (0.910,1.642) -- (0.929,1.651) -- (0.955,1.662) -- (0.974,1.670) -- (1.000,1.680) -- (1.025,1.689) -- (1.045,1.696) -- (1.071,1.704) -- (1.097,1.711) -- (1.122,1.718) -- (1.148,1.724) -- (1.174,1.729) -- (1.200,1.734) -- (1.232,1.739) -- (1.258,1.741) -- (1.290,1.744) -- (1.322,1.744) -- (1.354,1.743) -- (1.386,1.741) -- (1.412,1.737) -- (1.438,1.733) -- (1.464,1.726) -- (1.489,1.718) -- (1.515,1.709) -- (1.534,1.700) -- (1.556,1.689) -- (1.577,1.676) -- (1.595,1.663) -- (1.612,1.650) -- (1.631,1.632) -- (1.644,1.618) -- (1.660,1.599) -- (1.674,1.579) -- (1.686,1.560) -- (1.697,1.541) -- (1.708,1.516) -- (1.716,1.496) -- (1.725,1.470) -- (1.731,1.444) -- (1.736,1.418) -- (1.740,1.391) -- (1.743,1.361) -- (1.744,1.328) -- (1.744,1.296) -- (1.742,1.264) -- (1.739,1.238) -- (1.735,1.206) -- (1.731,1.180) -- (1.726,1.155) -- (1.720,1.129) -- (1.713,1.103) -- (1.706,1.077) -- (1.698,1.052) -- (1.689,1.026) -- (1.682,1.006) -- (1.672,0.981) -- (1.663,0.958) -- (1.654,0.936) -- (1.644,0.913) -- (1.633,0.891) -- (1.624,0.871) -- (1.612,0.846) -- (1.601,0.826) -- (1.591,0.807) -- (1.579,0.786) -- (1.567,0.764) -- (1.554,0.742) -- (1.542,0.723) -- (1.530,0.704) -- (1.518,0.685) -- (1.505,0.665) -- (1.492,0.646) -- (1.479,0.627) -- (1.465,0.607) -- (1.451,0.588) -- (1.438,0.570) -- (1.425,0.552) -- (1.412,0.536) -- (1.398,0.517) -- (1.382,0.498) -- (1.367,0.479) -- (1.354,0.463) -- (1.340,0.446) -- (1.323,0.427) -- (1.309,0.411) -- (1.295,0.395) -- (1.278,0.376) -- (1.264,0.360) -- (1.248,0.343) -- (1.232,0.326) -- (1.218,0.311) -- (1.200,0.292) -- (1.187,0.279) -- (1.168,0.260) -- (1.155,0.247) -- (1.135,0.227) -- (1.122,0.215) -- (1.103,0.196) -- (1.088,0.182) -- (1.071,0.166) -- (1.054,0.150) -- (1.039,0.136) -- (1.019,0.118) -- (1.005,0.105) -- (0.987,0.089) -- (0.970,0.073) -- (0.955,0.060) -- (0.936,0.042) -- (0.919,0.028) -- (0.903,0.014) -- (0.884,-0.003) -- (0.869,-0.017) -- (0.852,-0.032) -- (0.833,-0.049) -- (0.818,-0.062) -- (0.800,-0.078) -- (0.783,-0.094) -- (0.768,-0.108) -- (0.749,-0.125) -- (0.734,-0.139) -- (0.717,-0.156) -- (0.701,-0.172) -- (0.685,-0.188) -- (0.669,-0.204) -- (0.652,-0.221) -- (0.639,-0.236) -- (0.622,-0.255) -- (0.607,-0.272) -- (0.594,-0.288) -- (0.580,-0.307) -- (0.566,-0.326) -- (0.552,-0.345) -- (0.540,-0.365) -- (0.529,-0.384) -- (0.517,-0.405) -- (0.505,-0.429) -- (0.496,-0.448) -- (0.486,-0.474) -- (0.479,-0.497) -- (0.472,-0.519) -- (0.465,-0.545) -- (0.459,-0.577) -- (0.456,-0.603) -- (0.453,-0.635) -- (0.452,-0.661) -- (0.452,-0.693) -- (0.454,-0.719) -- (0.456,-0.751) -- (0.460,-0.777) -- (0.465,-0.809) -- (0.470,-0.835) -- (0.476,-0.861) -- (0.483,-0.886) -- (0.490,-0.912) -- (0.498,-0.937) -- (0.504,-0.957) -- (0.513,-0.983) -- (0.523,-1.009) -- (0.531,-1.028) -- (0.541,-1.054) -- (0.550,-1.073) -- (0.561,-1.099) -- (0.570,-1.118) -- (0.582,-1.142) -- (0.592,-1.163) -- (0.602,-1.182) -- (0.614,-1.205) -- (0.626,-1.227) -- (0.636,-1.247) -- (0.647,-1.266) -- (0.659,-1.286) -- (0.672,-1.308) -- (0.685,-1.330) -- (0.697,-1.350) -- (0.709,-1.369) -- (0.721,-1.388) -- (0.733,-1.408) -- (0.746,-1.427) -- (0.758,-1.446) -- (0.771,-1.466) -- (0.784,-1.485) -- (0.797,-1.504) -- (0.810,-1.524) -- (0.824,-1.543) -- (0.838,-1.562) -- (0.851,-1.582) -- (0.865,-1.600) -- (0.878,-1.618) -- (0.891,-1.636) -- (0.903,-1.653) -- (0.917,-1.672) -- (0.932,-1.691) -- (0.946,-1.710) -- (0.961,-1.730) -- (0.974,-1.747) -- (0.987,-1.763) -- (1.001,-1.781) -- (1.016,-1.800) -- (1.031,-1.820) -- (1.045,-1.837) -- (1.058,-1.853) -- (1.072,-1.871) -- (1.088,-1.891) -- (1.103,-1.909) -- (1.116,-1.925) -- (1.130,-1.942) -- (1.146,-1.961) -- (1.161,-1.980) -- (1.174,-1.995);
\draw[otNavy!30,line width=0.6pt] (-2.000,1.300) -- (-1.981,1.286) -- (-1.961,1.272) -- (-1.942,1.258) -- (-1.923,1.245) -- (-1.904,1.232) -- (-1.886,1.219) -- (-1.867,1.206) -- (-1.848,1.193) -- (-1.829,1.180) -- (-1.810,1.167) -- (-1.791,1.155) -- (-1.771,1.142) -- (-1.751,1.129) -- (-1.731,1.116) -- (-1.710,1.103) -- (-1.691,1.091) -- (-1.672,1.079) -- (-1.652,1.068) -- (-1.633,1.056) -- (-1.613,1.045) -- (-1.591,1.032) -- (-1.569,1.020) -- (-1.549,1.009) -- (-1.530,0.999) -- (-1.508,0.987) -- (-1.485,0.976) -- (-1.466,0.966) -- (-1.443,0.955) -- (-1.421,0.944) -- (-1.401,0.936) -- (-1.376,0.924) -- (-1.356,0.916) -- (-1.330,0.905) -- (-1.310,0.897) -- (-1.285,0.888) -- (-1.260,0.879) -- (-1.238,0.871) -- (-1.215,0.864) -- (-1.189,0.856) -- (-1.163,0.849) -- (-1.137,0.842) -- (-1.112,0.836) -- (-1.086,0.831) -- (-1.060,0.826) -- (-1.028,0.821) -- (-1.002,0.818) -- (-0.970,0.815) -- (-0.944,0.813) -- (-0.912,0.812) -- (-0.880,0.813) -- (-0.854,0.814) -- (-0.822,0.816) -- (-0.790,0.820) -- (-0.764,0.824) -- (-0.738,0.828) -- (-0.712,0.834) -- (-0.687,0.840) -- (-0.661,0.847) -- (-0.635,0.855) -- (-0.609,0.864) -- (-0.590,0.872) -- (-0.564,0.882) -- (-0.545,0.891) -- (-0.519,0.902) -- (-0.500,0.912) -- (-0.478,0.923) -- (-0.455,0.935) -- (-0.436,0.946) -- (-0.416,0.957) -- (-0.397,0.969) -- (-0.377,0.981) -- (-0.357,0.994) -- (-0.337,1.006) -- (-0.318,1.019) -- (-0.299,1.032) -- (-0.281,1.045) -- (-0.262,1.059) -- (-0.242,1.073) -- (-0.223,1.087) -- (-0.204,1.101) -- (-0.185,1.116) -- (-0.167,1.129) -- (-0.151,1.142) -- (-0.133,1.155) -- (-0.114,1.170) -- (-0.094,1.185) -- (-0.076,1.200) -- (-0.060,1.212) -- (-0.043,1.226) -- (-0.024,1.241) -- (-0.004,1.257) -- (0.013,1.270) -- (0.029,1.283) -- (0.047,1.297) -- (0.067,1.313) -- (0.086,1.328) -- (0.103,1.341) -- (0.119,1.354) -- (0.137,1.368) -- (0.157,1.383) -- (0.176,1.398) -- (0.194,1.412) -- (0.211,1.425) -- (0.228,1.438) -- (0.247,1.452) -- (0.266,1.466) -- (0.285,1.480) -- (0.305,1.494) -- (0.324,1.508) -- (0.342,1.521) -- (0.360,1.534) -- (0.379,1.547) -- (0.397,1.560) -- (0.416,1.573) -- (0.435,1.586) -- (0.455,1.599) -- (0.474,1.612) -- (0.494,1.624) -- (0.514,1.637) -- (0.535,1.650) -- (0.555,1.663) -- (0.575,1.675) -- (0.594,1.687) -- (0.614,1.698) -- (0.633,1.709) -- (0.654,1.721) -- (0.677,1.734) -- (0.697,1.745) -- (0.717,1.755) -- (0.738,1.766) -- (0.762,1.778) -- (0.781,1.788) -- (0.803,1.798) -- (0.826,1.809) -- (0.845,1.818) -- (0.871,1.829) -- (0.891,1.837) -- (0.916,1.848) -- (0.937,1.856) -- (0.961,1.865) -- (0.987,1.875) -- (1.008,1.882) -- (1.032,1.890) -- (1.058,1.898) -- (1.084,1.906) -- (1.109,1.913) -- (1.135,1.920) -- (1.161,1.926) -- (1.187,1.931) -- (1.212,1.936) -- (1.238,1.941) -- (1.270,1.946) -- (1.296,1.949) -- (1.328,1.952) -- (1.354,1.953) -- (1.386,1.954) -- (1.418,1.954) -- (1.448,1.953) -- (1.476,1.951) -- (1.509,1.947) -- (1.534,1.942) -- (1.560,1.937) -- (1.586,1.931) -- (1.612,1.923) -- (1.637,1.914) -- (1.657,1.906) -- (1.682,1.895) -- (1.702,1.885) -- (1.721,1.873) -- (1.740,1.861) -- (1.760,1.847) -- (1.779,1.832) -- (1.795,1.818) -- (1.811,1.802) -- (1.826,1.785) -- (1.842,1.766) -- (1.856,1.747) -- (1.869,1.728) -- (1.881,1.708) -- (1.891,1.689) -- (1.901,1.668) -- (1.912,1.644) -- (1.921,1.619) -- (1.927,1.599) -- (1.934,1.573) -- (1.940,1.547) -- (1.946,1.515) -- (1.949,1.489) -- (1.952,1.457) -- (1.954,1.431) -- (1.954,1.399) -- (1.954,1.367) -- (1.953,1.341) -- (1.950,1.309) -- (1.946,1.277) -- (1.943,1.251) -- (1.939,1.225) -- (1.933,1.197) -- (1.927,1.167) -- (1.921,1.142) -- (1.915,1.116) -- (1.908,1.090) -- (1.901,1.068) -- (1.894,1.045) -- (1.886,1.019) -- (1.877,0.994) -- (1.869,0.971) -- (1.861,0.948) -- (1.851,0.923) -- (1.843,0.903) -- (1.832,0.878) -- (1.823,0.858) -- (1.812,0.833) -- (1.803,0.813) -- (1.792,0.790) -- (1.781,0.768) -- (1.772,0.749) -- (1.760,0.726) -- (1.748,0.704) -- (1.738,0.685) -- (1.727,0.665) -- (1.715,0.643) -- (1.702,0.620) -- (1.690,0.601) -- (1.679,0.582) -- (1.667,0.562) -- (1.655,0.543) -- (1.643,0.524) -- (1.631,0.504) -- (1.618,0.484) -- (1.605,0.464) -- (1.592,0.445) -- (1.579,0.426) -- (1.567,0.407) -- (1.554,0.388) -- (1.540,0.369) -- (1.527,0.350) -- (1.513,0.330) -- (1.499,0.311) -- (1.485,0.292) -- (1.471,0.273) -- (1.457,0.254) -- (1.444,0.237) -- (1.431,0.220) -- (1.418,0.202) -- (1.403,0.182) -- (1.388,0.163) -- (1.373,0.144) -- (1.361,0.127) -- (1.348,0.111) -- (1.333,0.092) -- (1.318,0.073) -- (1.303,0.054) -- (1.290,0.038) -- (1.277,0.021) -- (1.262,0.002) -- (1.246,-0.017) -- (1.232,-0.035) -- (1.219,-0.052) -- (1.206,-0.069) -- (1.190,-0.088) -- (1.175,-0.107) -- (1.161,-0.126) -- (1.148,-0.142) -- (1.135,-0.159) -- (1.121,-0.178) -- (1.106,-0.197) -- (1.092,-0.217) -- (1.077,-0.236) -- (1.064,-0.254) -- (1.052,-0.272) -- (1.039,-0.290) -- (1.026,-0.309) -- (1.013,-0.328) -- (1.000,-0.347) -- (0.987,-0.367) -- (0.974,-0.388) -- (0.961,-0.409) -- (0.950,-0.429) -- (0.939,-0.448) -- (0.928,-0.468) -- (0.916,-0.491) -- (0.905,-0.513) -- (0.896,-0.532) -- (0.885,-0.558) -- (0.877,-0.577) -- (0.867,-0.603) -- (0.858,-0.627) -- (0.851,-0.648) -- (0.844,-0.674) -- (0.837,-0.700) -- (0.831,-0.725) -- (0.826,-0.751) -- (0.821,-0.783) -- (0.817,-0.809) -- (0.814,-0.841) -- (0.813,-0.867) -- (0.812,-0.899) -- (0.813,-0.931) -- (0.814,-0.957) -- (0.817,-0.989) -- (0.820,-1.017) -- (0.824,-1.047) -- (0.828,-1.073) -- (0.833,-1.099) -- (0.839,-1.124) -- (0.845,-1.150) -- (0.852,-1.176) -- (0.860,-1.202) -- (0.868,-1.227) -- (0.876,-1.253) -- (0.884,-1.275) -- (0.893,-1.298) -- (0.903,-1.324) -- (0.911,-1.343) -- (0.921,-1.369) -- (0.930,-1.388) -- (0.941,-1.414) -- (0.950,-1.433) -- (0.961,-1.456) -- (0.972,-1.479) -- (0.982,-1.498) -- (0.994,-1.520) -- (1.006,-1.543) -- (1.016,-1.562) -- (1.027,-1.582) -- (1.039,-1.602) -- (1.052,-1.624) -- (1.064,-1.646) -- (1.076,-1.665) -- (1.087,-1.685) -- (1.099,-1.704) -- (1.111,-1.723) -- (1.123,-1.742) -- (1.136,-1.762) -- (1.148,-1.781) -- (1.161,-1.800) -- (1.174,-1.820) -- (1.187,-1.839) -- (1.200,-1.858) -- (1.213,-1.878) -- (1.227,-1.897) -- (1.240,-1.916) -- (1.254,-1.936) -- (1.268,-1.955) -- (1.282,-1.974) -- (1.296,-1.994);

    \draw[otMuted,line width=0.5pt,dash pattern=on 2pt off 1.5pt] (-2.2,1.2) -- (1.2,-2.2);
    \draw[otOrange,line width=1.3pt] plot[domain=-0.58:0,samples=60]
      ({3*\x/(1+\x*\x*\x)},{3*\x*\x/(1+\x*\x*\x)});
    \draw[otOrange,line width=1.3pt] plot[domain=-0.58:0,samples=60]
      ({3*\x*\x/(1+\x*\x*\x)},{3*\x/(1+\x*\x*\x)});
    \draw[otOrange,line width=1.3pt] plot[domain=0:1,samples=80]
      ({3*\x/(1+\x*\x*\x)},{3*\x*\x/(1+\x*\x*\x)});
    \draw[otOrange,line width=1.3pt] plot[domain=0:1,samples=80]
      ({3*\x*\x/(1+\x*\x*\x)},{3*\x/(1+\x*\x*\x)});
  \end{scope}
  \draw[otMuted,line width=0.5pt,-{Stealth[length=1.5mm]}] (-2,0) -- (2.65,0);
  \draw[otMuted,line width=0.5pt,-{Stealth[length=1.5mm]}] (0,-2) -- (0,2.65);
  \fill[otGreenDark] (1,1) circle (1.8pt);
  \fill[otInk] (0,0) circle (1.6pt);
  \node[font=\scriptsize,text=otGreenDark,anchor=north west,fill=white,inner sep=1pt] at (1.04,0.96) {$(1,1)$};
  \node[font=\scriptsize,text=otMuted,anchor=north east] at (-0.92,-0.1) {$x+y+1=0$};
\end{tikzpicture}\\*[0.5mm]
\legende{Lignes de niveau de $f$ : le folium $f=0$ (orange), deux niveaux
négatifs (vert) entourant le minimum local $(1,1)$, deux niveaux positifs
(bleu). La droite pointillée est l'asymptote du folium.}
\end{minipage}\end{center}
\end{corrigeEval}

\begin{corrigeEval}[Problème]{7 points}{La suite $u_{n+1}=\sin u_n$}
\begin{questions}
  \item $\sin$ envoie $\bigl]0,\frac\pi2\bigr]$ dans $\intOF01\subset\bigl]0,\frac\pi2\bigr]$ :
    récurrence immédiate. Comme $\sin x<x$ pour $x>0$, $u_{n+1}<u_n$ : la
    suite décroît, minorée par $0$, donc converge vers
    $\ell\in\bigl[0,\frac\pi2\bigr]$, avec $\sin\ell=\ell$ par continuité ;
    d'où $\res{\ell=0}$. \pts{1 pt}
  \item Avec $\sin^{2}x=x^{2}-\frac{x^{4}}3+o(x^{4})$ :
    \[\frac{1}{\sin^{2}x}-\frac{1}{x^{2}}=\frac{x^{2}-\sin^{2}x}{x^{2}\sin^{2}x}
      =\frac{\frac{x^{4}}{3}+o(x^{4})}{x^{4}+o(x^{4})}\longrightarrow\frac13.\]
    Comme $u_n\to0$ avec $u_n>0$, par composition $\res{v_n\to\dfrac13}$.
    \pts{1,5 pt}
  \item Soit $\varepsilon>0$ et $N$ tel que $|a_k-\ell|\leqslant\varepsilon$ pour
    $k\geqslant N$. Pour $n>N$ :
    \[\Bigl|\frac1n\sum_{k=0}^{n-1}a_k-\ell\Bigr|\leqslant
      \underbrace{\frac1n\sum_{k=0}^{N-1}|a_k-\ell|}_{\to0\ (N\text{ fixé})}
      +\frac{n-N}{n}\,\varepsilon\leqslant2\varepsilon\]
    pour $n$ assez grand : \res{\text{le lemme de Cesàro}} est établi.
    \pts{1,5 pt}
  \item Par télescopage,
    $\frac1n\sum_{k=0}^{n-1}v_k=\frac1n\Bigl(\frac{1}{u_n^{2}}-\frac{1}{u_0^{2}}\Bigr)\to\frac13$,
    et $\frac{1}{nu_0^{2}}\to0$ : $\frac{1}{nu_n^{2}}\to\frac13$, soit
    $u_n^{2}\sim\frac3n$, et, $u_n$ étant positif,
    $\res{u_n\sim\sqrt{\dfrac3n}}$. \pts{1,5 pt}
  \item Par comparaison de séries à termes positifs avec les séries de
    Riemann : $u_n\sim\frac{\sqrt3}{n^{1/2}}$ et $u_n^{2}\sim\frac3n$ :
    $\res{\sum u_n\text{ et }\sum u_n^{2}\text{ divergent}}$ ;
    $u_n^{3}\sim\frac{3\sqrt3}{n^{3/2}}$ : $\res{\sum u_n^{3}\text{ converge}}$.
    \pts{1,5 pt}
\end{questions}
\begin{center}\begin{minipage}{\linewidth}\centering
\begin{tikzpicture}[x=2.6cm,y=2.6cm,line cap=round,line join=round]
  \draw[otMuted,line width=0.5pt,-{Stealth[length=1.5mm]}] (0,0) -- (1.75,0);
  \draw[otMuted,line width=0.5pt,-{Stealth[length=1.5mm]}] (0,0) -- (0,1.15);
  \draw[otMuted,line width=0.5pt] (0,0) -- (1.12,1.12);
  \draw[otNavy,line width=1.1pt] plot[domain=0:1.68,samples=80] (\x,{sin(deg(\x))});
  \draw[otGreenDark,line width=0.8pt] (1.5708,0) -- (1.5708,1) -- (1,1) -- (1,0.8415) -- (0.8415,0.8415)
    -- (0.8415,0.7456) -- (0.7456,0.7456) -- (0.7456,0.6784) -- (0.6784,0.6784) -- (0.6784,0.6276);
  \foreach \u/\l in {1.5708/0,1/1,0.8415/2}
    \node[font=\scriptsize,text=otGreenDark,anchor=north] at (\u,-0.02) {$u_{\l}$};
  \node[font=\scriptsize,text=otNavy,anchor=south] at (1.45,1.0) {$y=\sin x$};
  \node[font=\scriptsize,text=otMuted,anchor=south east] at (1.1,1.1) {$y=x$};
\end{tikzpicture}\\*[0.5mm]
\legende{Les itérés pour $u_0=\frac\pi2$ : la courbe de $\sin$ est tangente
à la droite $y=x$ en $0$, d'où une convergence lente, en $\sqrt{3/n}$ ; par
exemple $u_n\sqrt n\approx1{,}726$ pour $n=1000$.}
\end{minipage}\end{center}
\end{corrigeEval}
